\chapterepigraph{The world was to me a secret which I desired to divine.}{Mary
Shelley}{\textit{Frankenstein}}

\begingroup
\let\clearpage\relax
\let\cleardoublepage\relax
\chapter*{Introduction}
\endgroup
% \chapter* leaves \@nobreak globally true, and the \everypar that clears it is lost with
% the group; without this reset the first \paragraph below gets no space above it.
\global\csname @nobreakfalse\endcsname

A system of many interacting electrons is described completely by its wavefunction, and
the equation that fixes it is known exactly; the entire difficulty is solving it. For $N$
electrons in two dimensions the wavefunction is a function of $2N$ coordinates at once.
Stored on a grid of $M$ points per coordinate it would require $M^{2N}$ numbers, which for
twenty electrons and a modest ten points per axis is $10^{40}$. The competition that makes
this function hard to approximate is physical. Kinetic energy tends to delocalise the
particles, while the interactions and the external potential drive localisation and
correlation.

Established methods make one of two trades to remain tractable: they accept a cost that
grows combinatorially, or they decide in advance which correlations may be described. Full
configuration interaction makes the first. It is exact within a finite basis, but the
number of determinants it requires grows so quickly with the basis and with $N$ that a
modest two-dimensional dot of twenty electrons is already far beyond reach. Hartree--Fock
and coupled-cluster theory make the second. Hartree--Fock keeps only an average field, and
coupled cluster truncates the excitations it includes, which works well when correlation is weak and less well when it is strong and several
configurations contribute comparably. Fixed-node diffusion Monte Carlo makes the same kind
of decision through the nodes of its trial wavefunction. It provides accurate benchmarks
for continuum systems, but its accuracy is bounded by those nodes.

Variational Monte Carlo, the route this thesis takes, places the whole of that decision in
the trial wavefunction. The energy of any trial wavefunction bounds the ground-state energy
from above and can be written as an average of the local energy $\hat H\Psi/\Psi$ over
configurations distributed as $|\Psi|^2$. The trial state therefore never has to be
tabulated. It only has to be evaluated where the samples fall, and improving it becomes an
optimisation problem. Because $|\Psi|^2$ is positive, the method has no sign problem; it
accepts any trial wavefunction that can be evaluated; and its cost grows polynomially, as
$N^3$ to evaluate a determinant-based ansatz and one further power of $N$ for the Laplacian
in the local energy. Its accuracy is set by the trial wavefunction, and whether that
accuracy is reached depends on the optimiser and on the samples. In practice, then, much of
the many-body problem reduces to a narrower one: choosing an expressive, well-conditioned
ansatz, and optimising and sampling it reliably.

\paragraph{Quantum dots as a controlled testbed.}
Two-dimensional parabolic quantum dots are a clean and tunable place to study that
narrower problem. As the trap frequency $\omega$ decreases, the single-particle level
spacing collapses while Coulomb repulsion becomes dominant, and the electrons localise into
rings and shells, so that a single system spans the range from Fermi-liquid-like
behaviour to Wigner-molecule physics. In the lab frame, rotational symmetry and quantum
fluctuations smear these patterns into smooth densities; in a co-rotating frame the
crystalline structure reappears, with a topology that depends on particle number. For $N{=}6$ and
$N{=}12$, classical and semiclassical studies have identified preferred shellings such as
$(1,5)$ and $(3,9)$ and mapped parts of the crossover
\cite{Egger_1999,Kong_2002,schweigert1994spectralpropertieschargedparticles,manninen2007metalclustersquantumdots}.

Quantum dots suit this thesis for two reasons. They are genuinely non-trivial: Coulomb
singularities, strong correlation and a real crossover appear already at moderate $N$.
And they are simple enough that accurate DMC references exist in many regimes and that
detailed structural diagnostics are computationally feasible. For fully quantum dots,
quantitative topology-resolved diagnostics and a systematic picture of the Wigner regime
remain relatively sparse, and it is not well understood how modern neural ans\"atze
represent these states internally.

\paragraph{Neural wavefunctions.}
The narrower problem of choosing and optimising an ansatz is where machine learning
enters. As a tool for representing complex, high-dimensional functions, it has taken two
broad forms in physics. Neural networks act as
\emph{surrogates}, learning maps from parameters to observables or approximating the
solution operators of partial differential equations, as physics-informed neural networks
and operator learning do. Or they act as \emph{variational ans\"atze} for quantum states,
where the network outputs the amplitude of a many-body configuration and is trained by
variational Monte Carlo. This thesis belongs to the second tradition and uses the first:
a physics-informed residual objective serves as the pretraining stage of the variational
method and, in weak form, as the Markov-chain-free collocation route of
Chapter~\ref{ch:kernel}, a training method in its own right.

Representing a quantum state by a neural network and optimising it variationally was
introduced by Carleo and Troyer~\cite{Carleo2017-NQS}. For continuum fermions the idea has
since matured into highly accurate \emph{ab initio} solvers such as
FermiNet~\cite{Pfau2020-FermiNet} and PauliNet~\cite{Hermann2020-PauliNet}, and into
explicit neural \emph{backflow} transformations~\cite{LuoClark2019-NeuralBackflow}, the
direct ancestor of the coordinate backflow used here. Message passing has recently been
built into neural quantum states for extended fermion
systems~\cite{Pescia2024-MessagePassingNQS,KimEtAl2023-UltracoldFermiNQS}. Against that
backdrop the present work is deliberately narrow. It does not push a new architecture to
ever larger systems. It asks why particular architectural and training choices succeed or
fail on a problem where the answer can be checked, and what the trained state can then
say about the physics. Accuracy is the precondition for that, not its point.

\paragraph{Why a neural ansatz.}
Universal approximation establishes representational capacity, not trainability. The
theorems guarantee that a sufficiently large network can \emph{represent} any suitably
regular function to arbitrary accuracy; they say nothing about whether gradient-based
optimisation will \emph{find} that representation, nor at what sample cost, and it is the
second question that limits practice. The case for a neural ansatz therefore cannot rest
on expressiveness, and has to be made on cost. A variational neural ansatz makes neither of
the two trades above: it requires no enumeration of a determinant space, and it fixes no
prior truncation of which correlations may be described. It pays instead in optimisation
and sampling, and its capacity can still have to grow with $N$ or with the accuracy
demanded (Chapter~\ref{ch:kernel}). What it gains is the freedom to concentrate capacity
wherever the solution actually varies. Whether it does so is an empirical question, and
this thesis treats it as one.

The sampling cost has two parts. Every sample needs the local energy, and with it a
Laplacian of the network, however the samples are produced. In variational Monte Carlo the
samples also come from a Markov chain: Metropolis steps are sequential, must equilibrate,
and produce correlated configurations, and local moves mix slowly between well-separated
regions of configuration space. That cost is highest in the Wigner regime, where the
electrons localise into shell arrangements that the chain has to move between. Drawing
configurations instead from an analytic proposal and reweighting them removes the chain,
and with it the equilibration and the correlation between samples. It does not remove the
Laplacian, and its price is the variance of the importance weights.
Chapter~\ref{ch:kernel} maps where that trade holds.

The problem also has three properties that make a neural ansatz difficult to train, and each
shapes the ansatz used here. The Coulomb interaction diverges where two electrons meet, and
the local energy diverges with it unless the wavefunction cancels the singularity, so the
electron--electron cusp is built in analytically rather than learned. The kinetic energy is
a second derivative, which amplifies any error in the curvature of the network, so the
network's inputs are smooth, scaled to the trap and soft at short range. And the nodal surface,
where the determinant changes sign, is where a backflow receives its learning signal under a
strong-form residual objective, and also where the local energy of any imperfect
wavefunction diverges. In every variant tried in Appendix~\ref{app:catch22}, this kept the
strong form from training a backflow, which is why the collocation route uses a weak form. Beyond these,
the ansatz builds in antisymmetry and permutation symmetry and leaves the network to learn
the remaining correlation. On top of a Slater determinant of harmonic-oscillator orbitals,
two learnable objects carry that correlation: a Jastrow correlator, which reshapes the
amplitude, and a coordinate backflow, which reshapes the nodes.
Chapter~\ref{ch:trial_wavefunction} gives the construction.

\paragraph{Beyond the energy.}
A further choice concerns not \emph{what} structure is imposed but \emph{how} the network
represents the electrons. Correlation is relational: what one electron should do depends on
where the others are. Architectures that let particles exchange information are therefore
natural candidates for both the correlator and the backflow. Each can be built in two forms.
A separable network encodes per-particle and pairwise features and pools them in a single
pass, as a DeepSet does; a \emph{message-passing} network lets the particles exchange and
update information through explicit edge states before a quantity is read out. At strong
confinement, with everything else held fixed, the two correlators reach the same energy to
four figures. The energy therefore cannot show what the relational channel contributes, and
any difference between the architectures has to be sought in other quantities.

Such quantities are available because a trained ansatz can be inspected. Its physical structure is
imposed, as in any method, but the correlations it develops within those constraints are
not, and the trained network is a concrete object that can be examined directly, which the
walker population of diffusion Monte Carlo is not. A twenty-electron dot is a function of forty coordinates, far too
many to tabulate, but it does not follow that the solution itself is forty-dimensional in
any useful sense. How many directions does a given
parametrisation actually use to represent the solved state? That is a question about the
solution and its representation, not about the equation, and the answer depends on how the
state is parametrised (\S\ref{sec:theory-tangent-kernel}).

We make that question quantitative through the geometry of the wavefunction's
log-derivatives. Collecting the parameter sensitivities of a trained ansatz over an
ensemble of configurations into a covariance matrix gives, in the simplest reading, a
principal-component analysis of the state: the eigenvalue distribution counts how many
directions are doing genuine work. The physically appropriate version of that covariance
is the quantum geometric tensor, and the corresponding count is the participation ratio of
its spectrum. The same tensor is the metric a natural-gradient optimiser inverts, and it is
estimated from whatever samples define the loss, so it connects the architecture to the
optimiser and to the sampling measure. Chapter~\ref{ch:kernel} develops the instrument and
applies it to all three.

\paragraph{Questions.}
\label{sec:research-questions}
When a variational calculation stalls, or two models reach the same energy by different
routes, several explanations remain open: the representation may be wrong for the problem,
the optimisation may be failing, the samples that define the loss may no longer describe
the state, or the physics itself may have become harder. This thesis is organised as an
attempt to separate those explanations. Its first three questions follow them; the fourth
uses the resulting states for the physics they were meant to describe.
\begin{enumerate}
  \item \textbf{Architecture.} What does a message-passing network learn that a separable
        one cannot? We compare the two forms in the correlator and in the backflow, holding
        everything else fixed.
  \item \textbf{Optimiser.} When does natural-gradient (stochastic reconfiguration)
        preconditioning change the result, and when is it unnecessary or harmful?
  \item \textbf{Paradigm.} Removing the Markov chain from the gradient isolates the
        sampling explanation. Can an accurate wavefunction be trained by importance-sampled
        collocation instead, and how far into the correlated regime does that reach before
        the samples stop describing the state?
  \item \textbf{Physics.} Can the trained wavefunction serve as an instrument for
        Wigner-molecule physics (shell topologies, angular order, melting), and what does
        it show about the crossover?
\end{enumerate}

\noindent The results take these in a different order, because a state has to be shown
accurate enough to trust before either its physics or the solver that produced it can be
interpreted. Chapter~\ref{ch:results} establishes the energies and reads the physics from
the states; Chapter~\ref{ch:kernel} turns to the solver.

The questions are answered with three kinds of measurement: the tangent geometry above,
structural diagnostics of the trained state, which read shell topology and angular order
against a random-angle baseline and a held-out polygon shell model, and sampling
diagnostics for the Markov-chain-free route. All are post-processing on $|\Psi|^2$
samples, so the same analysis would apply unchanged to a classical or quantum-chemical
wavefunction. Chapter~\ref{ch:analysis} defines them.

\paragraph{Main contributions.}
\begin{enumerate}
  \item A compact Slater--Jastrow--backflow ansatz that agrees to within $0.05\%$ with the
        published fixed-node references for $N\in\{2,6,12,20\}$ at the confinements where we
        compare with them ($\omega=1$, $0.5$, $0.1$), and extends smoothly into the Wigner
        regime, where its energies are predictions.
  \item Evidence that message passing changes the correlator as a state, not only as a
        parametrisation. At equal energy a message-passing correlator reaches a lower
        local-energy variance and keeps its leading tangent directions inside a compact
        dictionary of collective operators; under the parametrisations used here it also
        concentrates its Fisher weight in about half as many effective directions as a
        separable one. A smaller conventional backflow collapses to a single displacement
        mode at weak confinement where a larger message-passing one does not.
  \item A hypothesis for when natural-gradient preconditioning earns its cost. Poor
        conditioning alone does not decide it, and that much is measured. The failures
        suggest that the metric must also be estimable from the batch.
  \item Markov-chain-free training by importance-sampled collocation, mapped to its limits:
        within about a tenth of a percent of DMC at strong confinement and within $0.35\%$ on
        every seed across the whole confinement range at $N{=}6$ with an adaptive proposal; a proposal
        built from the measured shell geometry that raises the effective sample size
        fifteen- to thirtyfold over a fixed one; and a cascade that carries converged states
        to $\omega=0.0035$, below which they no longer survive the step to a weaker trap.
  \item Topology-resolved diagnostics of the Wigner crossover, which quantify, shell by
        shell, how far radial stiffening runs ahead of angular order, complementing classical and quantum
        studies of parabolic
        dots~\cite{Egger_1999,Kong_2002,schweigert1994spectralpropertieschargedparticles,Filinov_2001}.
\end{enumerate}

\paragraph{Thesis structure.}
This introduction forms Part~I. Part~II sets out the physics and machine-learning background
(Chapters~\ref{sec:theory} and~\ref{ch:ml_theory}). Part~III gives the ansatz, the
optimisation pipeline and the diagnostics (Chapters~\ref{ch:trial_wavefunction}
to~\ref{ch:analysis}). Part~IV reports the results: Chapter~\ref{ch:results} the energies
and the Wigner-molecule structure, and Chapter~\ref{ch:kernel} the analysis of
architecture, optimisation and sampling together with the Markov-chain-free training
programme. Part~V discusses and concludes. The appendices in Part~VI develop the
conditioning theory, the collocation--backflow obstruction that motivated the weak-form
objective, the full diagnostic tables and the reproducibility record.
